\subsection{The original ring-switching evaluation reduction}\label{sec:dpr:protocol}

We state an algebraically explicit specialization of
Diamond--Posen's Construction~4.1
\cite{DiamondPosen2024}. Its large-field
evaluation protocol is represented by a black-box
committed-oracle interface, not silently replaced
by our Chapter~6 trace checker.

\begin{definition}[Underlying large-field evaluation interface]\label{def:dpr:large-pcs}
An \emph{underlying large-field evaluation protocol}
consists of algorithms
\[
 \mathsf{Commit}'(T'),\quad
 \mathsf{Prove}'(C',z,u;T'),\quad
 \mathsf{Verify}'(C',z,u,\pi').
\]
Here $T'$ is a multilinear polynomial over $L$,
$C'$ is the commitment to its Boolean coefficient
table, $z\in L^h$, and $u\in L$.
Perfect correctness means every honest opening
for $u=T'(z)$ is accepted. For a \emph{fixed,
bound} $T'$ we assume the following
evaluation soundness property: for any causal
adversary and any choice of the claimed
$(z,u)$ arising from the higher-level protocol,
\[
 \Pr[\mathsf{Verify}'\text{ accepts}\ \land\
 u\neq T'(z)]\leq\varepsilon_{\mathrm{eval}}.
\]
This is an assumption on the external protocol,
not a theorem of the tensor algebra.
\end{definition}

\begin{definition}[Ring-switching prover and verifier]\label{def:dpr:full-protocol}
The public instance is
$(C',r_{\rm lo},r_{\rm hi},y_{\rm pub})$;
the prover holds the Boolean table $w$
that packs into the committed $T'$.
The ring-switching reduction runs as follows.
\begin{enumerate}[label=\textup{(\roman*)}]
 \item \textbf{Commit:} The prover packs
 $w$ into the table of $T'$ and
 invokes $\mathsf{Commit}'(T')$ to obtain $C'$.
 The commitment and its bound table are fixed
 before the interactive evaluation.
 \item \textbf{Tensor partial message:}
 The prover forms
 $S=\varphi_1(T')(\varphi_0(r_{\rm hi}))$
 and sends an alleged value $S_0\in A$.
 \item \textbf{Original-claim check:}
 The verifier requires
 \[
 \boxed{y_{\rm pub}=C_q(S_0).}
 \]
 This step reads every tensor column
 but does not reveal the packed polynomial.
 \item \textbf{Ring sumcheck:}
 Run Definition~\ref{def:dpr:sumcheck}
 starting with $s_0=S_0$
 for the degree-two polynomial
 $H(X)=\varphi_1(T')(X)
       \operatorname{eq}_h(\varphi_0(r_{\rm hi}),X)$.
 For every $i<h$, the prover sends $g_i\in A[X]$
 of degree at most $2$; the verifier
 checks $s_i=g_i(0)+g_i(1)$, then samples
 $z_i\leftarrow L$ and puts
 $s_{i+1}=g_i(\varphi_1(z_i))$.
 \item \textbf{Reduction to the large field:}
 The prover sends an alleged
 $u\in L$ for $T'(z)$, where
 $z=(z_0,\ldots,z_{h-1})$,
 and the verifier requires
 \[
 \boxed{s_h=\varphi_1(u)\,
 \operatorname{eq}_h(\varphi_0(r_{\rm hi}),\varphi_1(z)).}
 \]
 \item \textbf{Delegated evaluation:}
 The parties execute the large-field
 evaluation protocol and accept only if
 $\mathsf{Verify}'(C',z,u,\pi')=1$.
\end{enumerate}
The verifier accepts precisely when
all the displayed checks and the
external proof pass.
\end{definition}

\begin{lemma}[All verifier quantities are typed]\label{lem:dpr:welltyped}
Every equality checked in
Definition~\ref{def:dpr:full-protocol}
is between elements of the same
ring or field, and the final
evaluation request to the black box
is over $L$, not over the tensor
algebra $A$.
\end{lemma}
\begin{proof}
The element $S_0$ and each
$s_i,g_i(0),g_i(1),
g_i(\varphi_1(z_i))$ belong to $A$.
The contraction $C_q(S_0)$ is in $L$,
matching $y_{\rm pub}\in L$.
The terminal expression is a
product of two elements of $A$,
matching $s_h\in A$.
Finally $u$ and $T'(z)$ belong to
$L$, and $C'$ is a commitment
to an $L$-coefficient table.
\end{proof}

\begin{lemma}[The public terminal equality factor is computable]\label{lem:dpr:terminal-public}
The element
$\operatorname{eq}_h(\varphi_0(r_{\rm hi}),\varphi_1(z))$
is computable by the verifier
from public values without
accessing the witness table.
\end{lemma}
\begin{proof}
All coordinates of $r_{\rm hi}$
and the freshly sampled $z$
are public at the final step.
Both embeddings $\varphi_0,\varphi_1$
and the equality polynomial are
specified by the public field
representation. Evaluating their
finite product requires no private
coefficient of $T'$.
\end{proof}

\begin{definition}[Committed-word timing invariant]\label{def:dpr:timing}
A run is \emph{causal and bound} if:
$T'$ is fixed by its commitment before
the evaluation interaction;
$S_0$ is sent before the sumcheck
challenges; each $g_i$ is fixed
before $z_i$; the final $u$ may
depend on the sampled $z$ but
must be verified against the
already bound $T'$; and the
external evaluation protocol satisfies
the assumed error bound of
Definition~\ref{def:dpr:large-pcs}.
\end{definition}

\begin{remark}[Protocol name and attribution]
This is the \emph{unmasked} reduction
of Diamond--Posen, specialized to
$K=\F_2$, extension $L/K$ of degree
$n=2^\ell$, and $h$ remaining variables.
It is distinct from the \emph{masked}
field-affine variant developed in
Chapter~7. In particular, the
original prover sends $S_0$
in the clear, whereas the masking
program aims to prevent disclosure
of its witness-dependent columns.
\end{remark}
